% Options for packages loaded elsewhere
\PassOptionsToPackage{unicode}{hyperref}
\PassOptionsToPackage{hyphens}{url}
\documentclass[
  spanish,
  12pt,
]{article}
\usepackage{xcolor}
\usepackage[margin=2.5cm]{geometry}
\usepackage{amsmath,amssymb}
\setcounter{secnumdepth}{5}
\usepackage{iftex}
\ifPDFTeX
  \usepackage[T1]{fontenc}
  \usepackage[utf8]{inputenc}
  \usepackage{textcomp} % provide euro and other symbols
\else % if luatex or xetex
  \usepackage{unicode-math} % this also loads fontspec
  \defaultfontfeatures{Scale=MatchLowercase}
  \defaultfontfeatures[\rmfamily]{Ligatures=TeX,Scale=1}
\fi
\usepackage{lmodern}
\ifPDFTeX\else
  % xetex/luatex font selection
\fi
% Use upquote if available, for straight quotes in verbatim environments
\IfFileExists{upquote.sty}{\usepackage{upquote}}{}
\IfFileExists{microtype.sty}{% use microtype if available
  \usepackage[]{microtype}
  \UseMicrotypeSet[protrusion]{basicmath} % disable protrusion for tt fonts
}{}
\makeatletter
\@ifundefined{KOMAClassName}{% if non-KOMA class
  \IfFileExists{parskip.sty}{%
    \usepackage{parskip}
  }{% else
    \setlength{\parindent}{0pt}
    \setlength{\parskip}{6pt plus 2pt minus 1pt}}
}{% if KOMA class
  \KOMAoptions{parskip=half}}
\makeatother
\usepackage{longtable,booktabs,array}
\newcounter{none} % for unnumbered tables
\usepackage{calc} % for calculating minipage widths
% Correct order of tables after \paragraph or \subparagraph
\usepackage{etoolbox}
\makeatletter
\patchcmd\longtable{\par}{\if@noskipsec\mbox{}\fi\par}{}{}
\makeatother
% Allow footnotes in longtable head/foot
\IfFileExists{footnotehyper.sty}{\usepackage{footnotehyper}}{\usepackage{footnote}}
\makesavenoteenv{longtable}
\usepackage{graphicx}
\makeatletter
\newsavebox\pandoc@box
\newcommand*\pandocbounded[1]{% scales image to fit in text height/width
  \sbox\pandoc@box{#1}%
  \Gscale@div\@tempa{\textheight}{\dimexpr\ht\pandoc@box+\dp\pandoc@box\relax}%
  \Gscale@div\@tempb{\linewidth}{\wd\pandoc@box}%
  \ifdim\@tempb\p@<\@tempa\p@\let\@tempa\@tempb\fi% select the smaller of both
  \ifdim\@tempa\p@<\p@\scalebox{\@tempa}{\usebox\pandoc@box}%
  \else\usebox{\pandoc@box}%
  \fi%
}
% Set default figure placement to htbp
\def\fps@figure{htbp}
\makeatother
% definitions for citeproc citations
\NewDocumentCommand\citeproctext{}{}
\NewDocumentCommand\citeproc{mm}{%
  \begingroup\def\citeproctext{#2}\cite{#1}\endgroup}
\makeatletter
 % allow citations to break across lines
 \let\@cite@ofmt\@firstofone
 % avoid brackets around text for \cite:
 \def\@biblabel#1{}
 \def\@cite#1#2{{#1\if@tempswa , #2\fi}}
\makeatother
\newlength{\cslhangindent}
\setlength{\cslhangindent}{1.5em}
\newlength{\csllabelwidth}
\setlength{\csllabelwidth}{3em}
\newenvironment{CSLReferences}[2] % #1 hanging-indent, #2 entry-spacing
 {\begin{list}{}{%
  \setlength{\itemindent}{0pt}
  \setlength{\leftmargin}{0pt}
  \setlength{\parsep}{0pt}
  % turn on hanging indent if param 1 is 1
  \ifodd #1
   \setlength{\leftmargin}{\cslhangindent}
   \setlength{\itemindent}{-1\cslhangindent}
  \fi
  % set entry spacing
  \setlength{\itemsep}{#2\baselineskip}}}
 {\end{list}}
\usepackage{calc}
\newcommand{\CSLBlock}[1]{\hfill\break\parbox[t]{\linewidth}{\strut\ignorespaces#1\strut}}
\newcommand{\CSLLeftMargin}[1]{\parbox[t]{\csllabelwidth}{\strut#1\strut}}
\newcommand{\CSLRightInline}[1]{\parbox[t]{\linewidth - \csllabelwidth}{\strut#1\strut}}
\newcommand{\CSLIndent}[1]{\hspace{\cslhangindent}#1}
\ifLuaTeX
\usepackage[bidi=basic,shorthands=off]{babel}
\else
\usepackage[bidi=default,shorthands=off]{babel}
\fi
\ifLuaTeX
  \usepackage{selnolig} % disable illegal ligatures
\fi
\setlength{\emergencystretch}{3em} % prevent overfull lines
\providecommand{\tightlist}{%
  \setlength{\itemsep}{0pt}\setlength{\parskip}{0pt}}
\usepackage{float}
\usepackage{array}
\usepackage{bookmark}
\IfFileExists{xurl.sty}{\usepackage{xurl}}{} % add URL line breaks if available
\urlstyle{same}
\hypersetup{
  pdflang={es},
  hidelinks,
  pdfcreator={LaTeX via pandoc}}

\title{\textcolor{blue}{Propuesta de investigación }}
\author{\textbf{Maria Paula Camargo}\footnote{mcamargori.edu.co},
Universidad Nacional de Colombia\\
\textbf{Karen Barrantes }\footnote{\href{mailto:kbarrantes@unal.edu.co}{\nolinkurl{kbarrantes@unal.edu.co}}},
Universidad Nacional de Colombia}
\date{}

\begin{document}
\maketitle
\begin{abstract}
En esta propuesta queremos estudiar si existe una asociación entre el
consumo diario de cafeína y la prevalencia autorreportada de enfermedad
cardiovascular (ECV) en adultos de Estados Unidos, usando los datos
públicos de la encuesta NHANES 2017-2018. La idea es caracterizar a la
población encuestada, describir cómo se comporta la prevalencia de ECV
según el nivel de consumo de cafeína, y evaluar esa relación con modelos
estadísticos que permitan controlar el efecto de la edad, el sexo y
otras variables sociodemográficas, para saber si lo que se observa a
simple vista se sostiene o si en realidad se explica por esos factores.
\end{abstract}

{
\setcounter{tocdepth}{3}
\tableofcontents
}
\section{Problema de investigación}\label{problema-de-investigaciuxf3n}

\subsection{Tema o título}\label{tema-o-tuxedtulo}

\textbf{Consumo de cafeína y enfermedad cardiovascular en adultos de
Estados Unidos en un análisis de la encuesta NHANES 2017-2018.}

\subsection{Planteamiento del
problema}\label{planteamiento-del-problema}

El café es una de las bebidas más consumidas en el mundo, y desde hace
años se discute qué tanto afecta a la salud del corazón. Parte de esa
preocupación viene de que la cafeína puede subir momentáneamente la
presión arterial y acelerar el ritmo cardíaco, lo que llevó durante
mucho tiempo a verla como un posible factor de riesgo cardiovascular.
Sin embargo, cuando se han hecho seguimientos a grandes grupos de
personas durante varios años, el panorama cambia bastante; un
metaanálisis que combinó 36 estudios de cohorte con más de 1,2 millones
de participantes encontró que el riesgo no aumenta de forma lineal con
el consumo, sino que el menor riesgo aparece en quienes toman entre 3 y
5 tazas al día, sin que los consumos altos muestren un riesgo mayor.
Otras revisiones han llegado a resultados parecidos, y algunos análisis
de cohortes conocidas como Framingham, ARIC y CHS incluso relacionan un
mayor consumo de café con menor riesgo de falla cardíaca.

A pesar de esa evidencia de estudios longitudinales, nos parece
pertinente mirar esta pregunta con una fuente de datos distinta; la
Encuesta Nacional de Salud y Nutrición de Estados Unidos (NHANES), que
es pública, representativa a nivel nacional y combina información
detallada sobre la dieta con el autorreporte de diagnósticos médicos. Un
punto que nos llama la atención del problema es que tanto el consumo de
cafeína como la presencia de ECV están asociados a la edad y al sexo,
así que cualquier relación que se vea ``a simple vista'' entre cafeína y
ECV podría deberse en realidad a estos dos factores, y no a un efecto
propio de la cafeína.

Por eso se creé que vale la pena revisar, con los datos de NHANES, si
existe una asociación entre la ingesta diaria de cafeína y la
prevalencia de enfermedad cardiovascular en adultos de Estados Unidos,
teniendo en cuenta al mismo tiempo la edad, el sexo y otras
características sociodemográficas que también podrían estar detrás de
esa prevalencia (véase el diccionario de variables en el Anexo 1). Con
esto buscamos ver si las diferencias que se observan inicialmente se
mantienen una vez se controlan estas características, o si más bien
desaparecen.

\subsection{Pregunta de
investigación}\label{pregunta-de-investigaciuxf3n}

\textbf{¿Existe asociación entre la ingesta diaria de cafeína y la
prevalencia autorreportada de enfermedad cardiovascular en adultos de
Estados Unidos (NHANES 2017-2018), una vez se tienen en cuenta la edad,
el sexo y otras características sociodemográficas?}

\section{Justificación}\label{justificaciuxf3n}

Estudiar la relación entre la cafeína y la enfermedad cardiovascular
parece relevante tanto para la salud pública como para la formación
estadística. La cafeína es una sustancia que consume prácticamente toda
la población adulta y que además es fácil de modificar en la dieta
diaria; por eso, tener evidencia clara sobre si se asocia (o no) con la
ECV puede aportar a las recomendaciones sobre qué tanto consumo se
considera seguro.

Desde lo \textbf{estadístico}, trabajar con una base secundaria como
NHANES permite poner en práctica muchos aspectos como, describir una
muestra, comparar grupos, evaluar si una relación observada se debe en
realidad a otras variables, y ajustar modelos que tengan en cuenta
varios factores al mismo tiempo. También da espacio para ir un poco más
allá de lo básico si el tiempo y los datos lo permiten, por ejemplo
mirando si la relación entre cafeína y ECV es realmente lineal o si se
comporta distinto en niveles muy altos de consumo.

Desde lo \textbf{práctico}, los resultados pueden ayudar a aclarar si el
consumo habitual de cafeína (hasta unos niveles razonables) debería
considerarse un factor de riesgo cardiovascular, o si lo que parece una
asociación en un análisis simple en realidad se explica por la edad y el
sexo de quienes consumen más cafeína.

Por último, este trabajo también tiene sentido como ejercicio de
formación ya que obliga a convertir una pregunta de salud pública en
variables medibles a partir de datos reales, plantear hipótesis que se
puedan poner a prueba y elegir herramientas estadísticas que tengan
sentido con el tipo de datos (transversales) y con las limitaciones
propias de un estudio observacional. En cuanto a viabilidad, los datos
de NHANES son de acceso público y ya e cuenta con un primer acercamiento
exploratorio a la base, lo que hace que el proyecto sea realizable.

\section{Objetivos}\label{objetivos}

\subsection{Objetivo general}\label{objetivo-general}

\textbf{Evaluar la asociación entre la ingesta diaria de cafeína y la
prevalencia autorreportada de enfermedad cardiovascular en adultos de
Estados Unidos, a partir de los datos de la encuesta NHANES 2017-2018,
considerando el efecto de la edad, el sexo y otras características
sociodemográficas.}

\subsection{Objetivos específicos}\label{objetivos-especuxedficos}

\begin{enumerate}
\def\labelenumi{\arabic{enumi}.}
\item
  \textbf{Caracterizar} a los adultos de la muestra según variables
  sociodemográficas (edad, sexo, raza, nivel educativo, razón
  ingreso/pobreza), consumo de cafeína y prevalencia de enfermedad
  cardiovascular.
\item
  \textbf{Describir y comparar} la prevalencia de enfermedad
  cardiovascular según categorías de consumo de cafeína (0, 1-100,
  101-200, 201-400 y más de 400 mg/día).
\item
  \textbf{Evaluar si la edad y el sexo explican} la relación observada
  entre cafeína y ECV, mediante un análisis por grupos etarios que
  permita ver si la tendencia se mantiene dentro de cada grupo.
\item
  \textbf{Estimar} la asociación entre el consumo de cafeína y la
  enfermedad cardiovascular mediante modelos de regresión logística,
  ajustando por edad, sexo, raza, nivel educativo, razón ingreso/pobreza
  y energía total consumida.
\item
  \textbf{Explorar si la relación entre cafeína y ECV es lineal} o si se
  comporta distinto en niveles altos de consumo, como un análisis
  complementario a los modelos de regresión de los objetivos anteriores.
\item
  \textbf{Revisar qué tan estables son los resultados}, comprobando si
  se mantienen al excluir del análisis los consumos de cafeína
  inusualmente altos y al mirar por separado cada tipo de evento
  cardiovascular (enfermedad coronaria, infarto y accidente
  cerebrovascular).
\end{enumerate}

\section{Marco conceptual}\label{marco-conceptual}

{[}Desarrolle aquí el marco conceptual.{]}

\section{Metodología}\label{metodologuxeda}

\subsection{Enfoque y diseño de
investigación}\label{enfoque-y-diseuxf1o-de-investigaciuxf3n}

La investigación se desarrolla bajo un enfoque cuantitativo, dado que se
utilizan variables medibles y procedimientos estadísticos para estudiar
la asociación entre la ingesta diaria de cafeína y la prevalencia
autorreportada de enfermedad cardiovascular en adultos de Estados
Unidos. Asimismo, se incorporan características sociodemográficas, como
la edad y el sexo, para evaluar la asociación ajustada por posibles
factores de confusión.

El estudio es de tipo observacional y analítico, puesto que no se
interviene sobre la ingesta de cafeína ni se asignan exposiciones a los
participantes, sino que se analizan los datos recolectados previamente
por la National Health and Nutrition Examination Survey (NHANES) durante
el período 2017--2018. Su propósito es evaluar la existencia y magnitud
de una asociación entre las variables de interés, considerando otras
características de los participantes.

El diseño es no experimental y transversal, dado que se emplea
información correspondiente a un único ciclo de la encuesta, sin
realizar seguimiento de los participantes a lo largo del tiempo. Aunque
NHANES tiene componentes de recolección en diferentes momentos, el
presente estudio utiliza los datos del ciclo 2017--2018 para examinar la
ingesta de cafeína y los antecedentes autorreportados de enfermedad
cardiovascular. En consecuencia, el análisis permite estudiar la
asociación entre las variables y la prevalencia de la condición, pero no
establecer una relación temporal ni concluir causalidad.

Esta elección metodológica es coherente con la pregunta de
investigación, ya que permite evaluar si la ingesta diaria de cafeína se
asocia con la prevalencia autorreportada de enfermedad cardiovascular en
adultos de Estados Unidos, una vez consideradas la edad, el sexo y otras
características sociodemográficas, mediante técnicas estadísticas de
análisis multivariable.

\subsection{Población, unidad de análisis y
muestra}\label{poblaciuxf3n-unidad-de-anuxe1lisis-y-muestra}

\subsubsection{Población objetivo.}\label{poblaciuxf3n-objetivo.}

La población objetivo de la National Health and Nutrition Examination
Survey (NHANES) 2017--2018 está constituida por la población civil no
institucionalizada residente en los Estados Unidos. La encuesta recoge
información sobre las características demográficas, socioeconómicas,
nutricionales y de salud de los participantes, mediante entrevistas,
exámenes físicos y pruebas de laboratorio, según el componente evaluado.

\subsubsection{Unidad de análisis.}\label{unidad-de-anuxe1lisis.}

La unidad de análisis corresponde a las personas participantes de la
encuesta que cumplen los criterios de inclusión definidos para el
estudio. La selección específica de la población analítica dependerá de
las variables consideradas y de la disponibilidad de información
completa para el análisis.

\subsubsection{Marco o fuente de
información.}\label{marco-o-fuente-de-informaciuxf3n.}

Se utilizarán datos secundarios de acceso público provenientes de NHANES
2017--2018, disponibles a través del National Center for Health
Statistics (NCHS), perteneciente a los Centers for Disease Control and
Prevention (CDC). Por tanto, no se realizará un nuevo proceso de
selección muestral, sino que se trabajará con los registros recolectados
por la encuesta original.

\subsubsection{Diseño muestral.}\label{diseuxf1o-muestral.}

NHANES emplea un diseño de muestreo probabilístico, estratificado, por
conglomerados y multietápico, cuyo propósito es obtener estimaciones
representativas de la población civil no institucionalizada de los
Estados Unidos. El proceso de selección comprende cuatro etapas
principales: (i) selección de unidades primarias de muestreo,
constituidas por condados o grupos de condados; (ii) selección de
segmentos geográficos dentro de dichas unidades, definidos a partir de
bloques censales o agrupaciones de estos; (iii) selección de viviendas u
otras unidades residenciales dentro de los segmentos elegidos; y (iv)
selección de personas elegibles dentro de los hogares participantes.

El diseño contempla probabilidades desiguales de selección y
sobremuestreo de determinados grupos poblacionales para mejorar la
precisión de las estimaciones en subpoblaciones de interés. En
consecuencia, el análisis estadístico deberá considerar las
ponderaciones muestrales y las variables de estratificación y
conglomeración correspondientes al componente de NHANES utilizado.

\textbf{Criterios de inclusión y exclusión.} Se incluirán los
participantes pertenecientes al ciclo 2017--2018 que cumplan los
criterios de elegibilidad establecidos para la población de interés y
dispongan de información válida para las variables requeridas. Se
excluirán del análisis los registros que no cumplan dichos criterios o
que carezcan de información indispensable para las variables
principales, de acuerdo con la estrategia definida para el manejo de
datos faltantes. Estos criterios se precisarán según la pregunta de
investigación y las variables seleccionadas.

\textbf{Tamaño de muestra.} No se calculará un nuevo tamaño de muestra,
dado que se utilizará una base de datos existente. El tamaño de la
muestra analítica estará determinado por los participantes del ciclo
2017--2018 que cumplan los criterios de inclusión y cuenten con la
información necesaria para el estudio.

\subsection{Variables e indicadores}\label{variables-e-indicadores}

\begin{table}[H]
\centering
\small
\caption{Definición y caracterización de las variables del estudio}
\label{tab:variables}

\renewcommand{\arraystretch}{1.25}
\setlength{\tabcolsep}{3pt}

\begin{tabular}{
|>{\raggedright\arraybackslash}p{0.14\textwidth}
|>{\raggedright\arraybackslash}p{0.24\textwidth}
|>{\raggedright\arraybackslash}p{0.25\textwidth}
|>{\raggedright\arraybackslash}p{0.16\textwidth}
|>{\raggedright\arraybackslash}p{0.11\textwidth}|
}
\hline

\textbf{Variable} &
\textbf{Papel en el estudio} &
\textbf{Definición operacional} &
\textbf{Tipo/escala} &
\textbf{Fuente} \\
\hline

[ecv] &
[Respuesta] &
[Variable derivada de *acv, coronaria* e *infarto* (ver anexo 1), Se codifica como 1 si la persona reporta haber recibido al menos uno de los tres diagnósticos y como 0 si reporta no haber recibido ninguno.] &
[Cualitativa, nominal, binaria] &
[Derivada] \\
\hline

[cafeina\_mg] &
[explicativa] &
[Cafeína total
ingerida en el día
(recordatorio 24h), mg/día, "0 mg", "1-100 mg", "101-200 mg", "201-400 mg", ">400 mg"] &
[Cualitativa, nominal] &
[DR1TCAFF] \\
\hline

[sexo] &
[control] &
[Sexo, codificado como Hombre/Mujer] &
[Cualitativa] &
[RIAGENDR] \\
\hline

[raza] &
[control] &
[Raza y origen hispano, codificado como: Negro no hispano, Asiatico no hispano, Blanco no hispano, Otra Raza / multirracial, Mexicano-americano, otro hispano.] &
[Cualitativa, nominal] &
[RIDRETH3] \\
\hline

[educacion] &
[control] &
[Nivel educativo, codificado como: "Menos de 9no grado, 9-11 grado (sin terminar), Bachiller / GED, Universidad incompleta / tecnico,  Universitario o mas] &
[Cualitativa, ordinal] &
[DMDEDUC2] \\
\hline

[pir] &
[control] &
[Razón ingreso familiar / línea de pobreza] &
[Cuantitativa] &
[INDFMPIR] \\
\hline

[kcal] &
[control] &
[Energía total ingerida en el día, codificada como kcal/día] &
[Cuantitativa] &
[DR1TKCAL] \\
\hline

\end{tabular}
\end{table}

\subsection{Fuentes de información e
instrumentos}\label{fuentes-de-informaciuxf3n-e-instrumentos}

\subsubsection{Fuentes}\label{fuentes}

Los datos utilizados en el estudio provienen de la \emph{National Health
and Nutrition Examination Survey} (NHANES), correspondiente al ciclo
2017--2018. Esta encuesta es desarrollada por el \emph{National Center
for Health Statistics} (NCHS), entidad perteneciente a los \emph{Centers
for Disease Control and Prevention} (CDC) de los Estados Unidos. Su
objetivo es caracterizar el estado de salud y nutrición de la población
mediante información demográfica, socioeconómica, de salud y de
comportamiento, así como mediciones físicas y pruebas de laboratorio.

NHANES utiliza un diseño de muestreo probabilístico, estratificado, por
conglomerados y multietápico, dirigido a la población civil no
institucionalizada residente en los Estados Unidos. La información se
obtiene mediante entrevistas domiciliarias y evaluaciones realizadas en
centros de examen móviles, de acuerdo con los componentes de la
encuesta.

Para este estudio se emplearán datos secundarios de acceso público,
disponibles en el portal oficial de NHANES. En particular, se utilizará
información del cuestionario de condiciones médicas (\emph{Medical
Conditions Questionnaire}), que incluye antecedentes de enfermedades
reportados por los participantes. A partir de las variables
correspondientes a accidente cerebrovascular, enfermedad coronaria e
infarto de miocardio, se construirá una variable derivada denominada
\texttt{ecv}, que identificará a las personas que reporten al menos uno
de estos antecedentes.

Dado que se utilizará una base de datos existente, no se realizará
recolección primaria de información. El procesamiento comprenderá la
selección de las variables pertinentes, la revisión de los códigos de
respuesta y de los valores faltantes, la construcción de la variable
derivada y la aplicación de los criterios de inclusión y exclusión
definidos para el análisis. Asimismo, se tendrán en cuenta las
características del diseño muestral de NHANES en los análisis
estadísticos, cuando corresponda.

\subsubsection{Instrumentos}\label{instrumentos}

{[}Describa aquí las fuentes de información y los instrumentos.{]}

\subsection{Procedimiento}\label{procedimiento}

\subsubsection{Integración de las bases de
datos}\label{integraciuxf3n-de-las-bases-de-datos}

Para el desarrollo del análisis se utilizaron tres bases de datos de la
encuesta NHANES, correspondientes al ciclo 2017--marzo de 2020
(\emph{pre-pandemic}), identificadas mediante los archivos
\texttt{DEMO\_J.xpt}, \texttt{DR1TOT\_J.xpt} y \texttt{MCQ\_J.xpt}. De
cada base se seleccionaron únicamente las variables de interés para el
estudio.

\begin{itemize}
\tightlist
\item
  \textbf{DEMO\_J (Demographics):} contiene información sociodemográfica
  de los participantes, necesaria para caracterizar la población de
  estudio.
\item
  \textbf{DR1TOT\_J (Dietary Interview -- Total Nutrient Intakes, First
  Day):} contiene información sobre la ingesta dietética total y el
  consumo de nutrientes durante el primer día de recordatorio
  alimentario. De esta base se extrajeron las variables dietéticas
  requeridas para el análisis.
\item
  \textbf{MCQ\_J (Medical Conditions):} contiene información relacionada
  con las condiciones médicas reportadas por los participantes. Se
  seleccionaron las variables necesarias para identificar las
  condiciones de salud consideradas en el estudio.
\end{itemize}

Posteriormente, las tres bases se integraron mediante el identificador
único de cada participante (\texttt{SEQN}), utilizando un \emph{inner
join}. Este procedimiento permitió conservar únicamente los
participantes presentes en las tres fuentes de información y reunir sus
características sociodemográficas, dietéticas y de salud en una sola
base de datos.

Como resultado de este proceso de integración, se obtuvo una base de
datos consolidada con \textbf{8.366 registros}, cada uno correspondiente
a un participante que contaba con información en las tres bases
originales. Esta base constituyó el punto de partida para los procesos
posteriores de depuración, selección de variables y definición de la
muestra analítica.

\subsubsection{Selección de la población
estudio}\label{selecciuxf3n-de-la-poblaciuxf3n-estudio}

Se restringió la población analítica a participantes de 20 años o más,
debido a que las preguntas del cuestionario de condiciones médicas
utilizadas para construir la variable de enfermedad cardiovascular se
dirigen a este grupo etario. Este criterio permite mantener la
coherencia entre la población de estudio, la disponibilidad de la
variable de respuesta y la pregunta de investigación.

Después de aplicar el criterio de inclusión de participantes de 20 años
o más, la base de datos pasó de 8.366 a 5.265 registros.

\subsubsection{Tratamiento de datos faltantes y estrategia de
análisis}\label{tratamiento-de-datos-faltantes-y-estrategia-de-anuxe1lisis}

Posteriormente, se evaluó la presencia de valores faltantes en las
variables de ingesta diaria de cafeína (\texttt{DR1TCAFF}), ingesta
energética total (\texttt{DR1TKCAL}) y razón entre el ingreso familiar y
el umbral federal de pobreza (\texttt{INDFMPIR}). Se identificaron 523
participantes con información faltante en las variables dietarias y 697
con información faltante en el PIR. La revisión de los patrones de
ausencia mostró que los faltantes de cafeína y energía coincidían en los
mismos participantes, mientras que los correspondientes al PIR
presentaban una superposición parcial con los anteriores.

En total, 1.081 de los 5.265 participantes que cumplían el criterio de
edad presentaron información faltante en al menos una de estas
variables, lo que representa el 20,5 \% de la población elegible.
Adicionalmente, durante la recodificación del nivel educativo
(\texttt{DMDEDUC2}), los códigos correspondientes a respuestas no
concluyentes se trataron como valores faltantes, de acuerdo con la
codificación original de la encuesta. Asimismo, se identificó un
registro con información insuficiente para determinar el desenlace de
enfermedad cardiovascular autorreportada (\texttt{ecv}). Una vez
aplicados los criterios de inclusión y disponibilidad de información, se
estableció una muestra analítica preliminar de casos completos de 4.180
participantes adultos, cifra que deberá corresponder al resultado
verificado de la depuración final.

Para el análisis principal de la asociación entre la ingesta de cafeína
y los antecedentes autorreportados de enfermedad cardiovascular, se
utilizará una estrategia de análisis de casos completos
(\emph{complete-case analysis}). En consecuencia, la regresión logística
multivariada incluirá únicamente a los participantes con información
válida y observada en el desenlace y en todas las variables explicativas
y de ajuste incorporadas en el modelo. Esta decisión permite definir una
muestra analítica consistente para la estimación de los \emph{odds
ratios} ajustados. No obstante, de acuerdo con Little y Rubin
(\citeproc{ref-little2020statistical}{2020}) (2019, cap. 3), la validez
de este procedimiento depende de los supuestos relativos al mecanismo de
ausencia y a la estructura del modelo; por tanto, no se asumirá que las
estimaciones están libres de sesgo únicamente porque los faltantes se
concentren en las covariables.

Como parte del diagnóstico de los datos faltantes, se compararon
descriptivamente algunas características de los participantes con
información completa e incompleta. La edad media fue de 51,2 años entre
los casos completos y de 52,2 años entre los incompletos; la proporción
de mujeres fue del 51,4 \% y del 53,2 \%, respectivamente. La
prevalencia descriptiva de enfermedad cardiovascular autorreportada fue
del 10,4 \% en el grupo completo y del 9,0 \% en el incompleto. Estas
diferencias permiten caracterizar ambos grupos, pero no son suficientes
para determinar el mecanismo de ausencia ni para descartar posibles
sesgos derivados de la exclusión de registros.

Como estrategia complementaria, se propone realizar un análisis de
sensibilidad mediante imputación múltiple, siempre que la evaluación de
los patrones de ausencia y la información disponible permitan justificar
este procedimiento. Para ello, se generarán varias bases de datos
imputadas mediante un modelo que considere el desenlace, las variables
explicativas y de ajuste, y otras variables auxiliares pertinentes. La
regresión logística se estimará en cada base imputada y los resultados
se combinarán mediante las reglas de Rubin
(\citeproc{ref-little2020statistical}{Little y Rubin 2020}), teniendo en
cuenta las particularidades del diseño muestral de NHANES.
Posteriormente, se compararán los \emph{odds ratios} ajustados y sus
intervalos de confianza con los obtenidos en el análisis principal de
casos completos, con el propósito de evaluar la robustez de las
conclusiones frente al tratamiento de los datos faltantes.

Finalmente, se reconoce que el análisis de casos completos puede reducir
la precisión de las estimaciones y modificar la composición de la
muestra analítica. Por ello, la comparación con los resultados de la
imputación múltiple permitirá examinar si las conclusiones sustantivas
dependen de la estrategia utilizada para tratar los faltantes. La
interpretación definitiva de las estimaciones considerará, además, los
pesos, estratos y unidades primarias de muestreo correspondientes al
diseño complejo de NHANES.

\subsubsection{Valores atípicos y
consistencia}\label{valores-atuxedpicos-y-consistencia}

En cuanto a la revisión de valores atípicos, se examinarán la
distribución, los rangos y los valores extremos de las variables
cuantitativas, especialmente la ingesta de cafeína, la ingesta
energética y la edad. Se utilizarán resúmenes estadísticos y
representaciones gráficas, como histogramas y diagramas de caja, para
identificar observaciones potencialmente inusuales. Los valores extremos
no se eliminarán automáticamente, dado que pueden corresponder a
consumos reales; cualquier corrección o exclusión requerirá evidencia de
error de registro, inconsistencia con la documentación de la encuesta o
una justificación metodológica explícita.

La consistencia de los datos se verificará mediante la revisión de los
códigos válidos y las categorías de las variables, la detección de
registros duplicados, la comprobación de los rangos esperados y la
identificación de posibles contradicciones entre las respuestas de las
condiciones médicas. También se comprobará que las variables derivadas
sean coherentes con sus componentes originales y que los valores
faltantes se propaguen adecuadamente cuando no exista información
suficiente para realizar una clasificación.

\subsubsection{Variables derivadas}\label{variables-derivadas}

Finalmente, se construirán las variables derivadas necesarias para el
análisis. La variable de enfermedad cardiovascular (ecv) se definirá
como una variable binaria que toma el valor de uno cuando el
participante reporta al menos una de las siguientes condiciones:
enfermedad coronaria, infarto de miocardio o accidente cerebrovascular;
tomará el valor de cero cuando las tres respuestas indiquen ausencia de
dichas condiciones. Si las respuestas disponibles no permiten establecer
ninguna de estas dos situaciones, el desenlace se conservará como
faltante. Adicionalmente, la ingesta diaria de cafeína se agrupará en
las categorías de 0 mg, 1-100 mg, 101-200 mg, 201-400 mg y más de 400
mg. La edad se clasificará en los grupos de 20-39, 40-59 y 60 años o
más. Estas categorizaciones permitirán describir la distribución de la
exposición y de las características de los participantes. Se conservarán
también las variables cuantitativas originales para evitar perder
información cuando se evalúe la ingesta de cafeína como variable
continua.

\subsection{Plan preliminar de análisis
estadístico}\label{plan-preliminar-de-anuxe1lisis-estaduxedstico}

{\def\LTcaptype{none} % do not increment counter
\begin{longtable}[]{@{}
  >{\raggedright\arraybackslash}p{(\linewidth - 4\tabcolsep) * \real{0.2603}}
  >{\raggedright\arraybackslash}p{(\linewidth - 4\tabcolsep) * \real{0.3014}}
  >{\raggedright\arraybackslash}p{(\linewidth - 4\tabcolsep) * \real{0.4384}}@{}}
\toprule\noalign{}
\begin{minipage}[b]{\linewidth}\raggedright
Objetivo específico
\end{minipage} & \begin{minipage}[b]{\linewidth}\raggedright
Variables involucradas
\end{minipage} & \begin{minipage}[b]{\linewidth}\raggedright
Procedimiento estadístico propuesto
\end{minipage} \\
\midrule\noalign{}
\endhead
\bottomrule\noalign{}
\endlastfoot
1 & Edad, sexo, raza/etnia, nivel educativo, razón ingreso/pobreza ,
consumo de cafeína y ECV. & Se realizará un análisis descriptivo de las
características de la población analítica. Para las variables
categóricas se estimarán frecuencias y porcentajes; para las
cuantitativas se calcularán medidas de tendencia central y dispersión,
como media, desviación estándar, mediana y rango intercuartílico, según
la distribución observada. Se estimarán la prevalencia de ECV y la
distribución del consumo de cafeína. Se incorporarán los pesos y el
diseño muestral de NHANES para obtener estimaciones representativas de
la población objetivo. \\
2 & ECV (variable dicotómica) y consumo de cafeína & Se estimará la
prevalencia de ECV en cada categoría de consumo, acompañada de
intervalos de confianza del 95 \%. Para evaluar si existen diferencias
entre categorías se empleará una prueba de asociación apropiada para
encuestas complejas, como la prueba de Rao--Scott, que ajusta la prueba
de independencia por el diseño muestral. Esto permitirá identificar
diferencias globales en la prevalencia, sin asumir que estas se deben
exclusivamente al consumo de cafeína. \\
3 & Consumo de cafeína, ECV, edad y sexo. & Se comparará la asociación
observada inicialmente entre cafeína y ECV con la obtenida al considerar
la edad y el sexo. Además, se estimarán las prevalencias de ECV por
categoría de cafeína dentro de los grupos etarios definidos (20--39,
40--59 y 60 años o más), para examinar si la tendencia observada se
mantiene entre edades. Se considerará la inclusión de términos de
interacción entre cafeína y edad, y entre cafeína y sexo, en modelos de
regresión logística, para evaluar si la asociación varía según estas
características. Así se podrá distinguir entre una relación que cambia
al controlar estas variables y otra cuya magnitud difiere entre
grupos. \\
4 & ECV, consumo de cafeína, edad, sexo, raza/etnia, nivel educativo,
razón ingreso/pobreza y energía total consumida. & Se ajustará un modelo
de regresión logística para encuestas complejas, con ECV como variable
dependiente y consumo de cafeína como variable explicativa principal. Se
incluirán las covariables especificadas para controlar posibles factores
de confusión. Se reportarán razones de momios (odds ratios, OR)
ajustadas e intervalos de confianza del 95 \%. El modelo permitirá
evaluar la asociación entre el consumo de cafeína y la probabilidad de
reportar antecedentes de ECV, condicionada a las covariables incluidas.
La interpretación se hará como asociación, no como efecto causal. \\
5 & Consumo de cafeína en mg/día y ECV, con las mismas covariables de
ajuste del objetivo 4. & Se comparará el modelo que incorpora cafeína
como variable continua lineal con un modelo que permita una relación no
lineal, por ejemplo, mediante splines cúbicos restringidos. Esto
permitirá examinar si el cambio en la asociación es aproximadamente
constante a lo largo del consumo o si existen curvaturas o variaciones
en niveles altos. La forma de la relación se evaluará mediante pruebas
de no linealidad y la representación gráfica de las asociaciones
estimadas con sus intervalos de confianza. Este análisis complementará
el modelo principal sin depender exclusivamente de categorías
arbitrarias de consumo. \\
6 & Consumo de cafeína, ECV global y cada evento por separado:
enfermedad coronaria (MCQ160C), infarto (MCQ160E) y accidente
cerebrovascular (MCQ160F). & Se realizarán análisis de sensibilidad.
Primero, se repetirá el modelo principal excluyendo consumos
inusualmente altos, definidos mediante un criterio explícito y
documentado, para determinar si las estimaciones dependen de
observaciones extremas. Segundo, se ajustarán modelos separados para
cada evento cardiovascular, conservando las covariables y el diseño
muestral. Se compararán las OR ajustadas, sus intervalos de confianza y
la dirección y magnitud de las asociaciones con los resultados
principales. Estos análisis permitirán valorar la robustez de los
hallazgos y si la asociación difiere según el evento estudiado. \\
\end{longtable}
}

\subsection{Consideraciones éticas y manejo de la
información}\label{consideraciones-uxe9ticas-y-manejo-de-la-informaciuxf3n}

El estudio utiliza información secundaria proveniente de la
\emph{National Health and Nutrition Examination Survey} (NHANES)
correspondiente al periodo 2017-2018, realizada por el National Center
for Health Statistics (NCHS) de los Centers for Disease Control and
Prevention (CDC) de Estados Unidos. Por tanto, no implica el
reclutamiento de participantes ni la recolección directa de datos por
parte del equipo investigador.

\textbf{Consentimiento informado.} El consentimiento informado
corresponde al proceso de recolección original de NHANES, en el cual se
obtiene la autorización de los participantes para su participación en la
encuesta y los procedimientos asociados. En este estudio no se
solicitará un nuevo consentimiento, dado que se analizarán datos
previamente recolectados y disponibles públicamente.

\textbf{Privacidad, confidencialidad y anonimización.} Se utilizará la
base de datos pública de NHANES, en la que se han aplicado medidas de
protección de la confidencialidad y no se incluyen identificadores
personales directos que permitan reconocer a los participantes. El
análisis se realizará sobre registros individuales desidentificados y
los resultados se presentarán de forma agregada, evitando cualquier
intento de identificación o divulgación de información que pueda
comprometer la privacidad de las personas.

\textbf{Manejo seguro de la información.} Los archivos descargados se
almacenarán en un entorno de trabajo con acceso controlado y se
utilizarán exclusivamente para los fines académicos de la investigación.
Durante la integración, depuración y análisis de las bases de datos se
conservarán únicamente las variables necesarias para cumplir los
objetivos del estudio, sin intentar vincular los registros con fuentes
externas para identificar participantes.

\textbf{Uso responsable y limitaciones.} La información se analizará de
acuerdo con los lineamientos de uso de datos públicos de NHANES y se
reconocerá la fuente original. Las conclusiones se formularán
considerando las limitaciones del diseño observacional y transversal de
la encuesta, por lo que las asociaciones encontradas entre el consumo de
cafeína y la enfermedad cardiovascular no se interpretarán como
evidencia de causalidad.

Dado que se trata de un análisis secundario de datos públicos
desidentificados, no se prevé contacto con participantes ni recolección
de información identificable. Los requisitos institucionales aplicables
a la revisión ética de este tipo de investigación deberán verificarse
según las normas de la institución académica correspondiente.

\section{Cronograma}\label{cronograma}

{\def\LTcaptype{none} % do not increment counter
\begin{longtable}[]{@{}
  >{\raggedright\arraybackslash}p{(\linewidth - 8\tabcolsep) * \real{0.4359}}
  >{\centering\arraybackslash}p{(\linewidth - 8\tabcolsep) * \real{0.1410}}
  >{\centering\arraybackslash}p{(\linewidth - 8\tabcolsep) * \real{0.1410}}
  >{\centering\arraybackslash}p{(\linewidth - 8\tabcolsep) * \real{0.1410}}
  >{\centering\arraybackslash}p{(\linewidth - 8\tabcolsep) * \real{0.1410}}@{}}
\toprule\noalign{}
\begin{minipage}[b]{\linewidth}\raggedright
Actividad
\end{minipage} & \begin{minipage}[b]{\linewidth}\centering
Período 1
\end{minipage} & \begin{minipage}[b]{\linewidth}\centering
Período 2
\end{minipage} & \begin{minipage}[b]{\linewidth}\centering
Período 3
\end{minipage} & \begin{minipage}[b]{\linewidth}\centering
Período 4
\end{minipage} \\
\midrule\noalign{}
\endhead
\bottomrule\noalign{}
\endlastfoot
Revisión bibliográfica & X & X & & \\
Obtención y preparación de datos & & X & X & \\
Análisis estadístico & & & X & X \\
Elaboración del informe final & & & & X \\
\end{longtable}
}

\section{Anexos}\label{anexos}

Incluya, cuando corresponda:

\begin{itemize}
\tightlist
\item
  Tablas o gráficas adicionales.\\
\item
  Archivos de código R utilizados.\\
\item
  Encuestas, modelos o instrumentos.\\
\item
  Diccionario de variables o ficha técnica de la base de datos.\\
\item
  Discusiones teóricas sobre modelos o metodologías.
\end{itemize}

\section{Anexos}\label{anexos-1}

\textbf{Anexo 1. Diccionario de variables de la base analítica.}

{\def\LTcaptype{none} % do not increment counter
\begin{longtable}[]{@{}
  >{\raggedright\arraybackslash}p{(\linewidth - 6\tabcolsep) * \real{0.2500}}
  >{\raggedright\arraybackslash}p{(\linewidth - 6\tabcolsep) * \real{0.2500}}
  >{\raggedright\arraybackslash}p{(\linewidth - 6\tabcolsep) * \real{0.2500}}
  >{\raggedright\arraybackslash}p{(\linewidth - 6\tabcolsep) * \real{0.2500}}@{}}
\toprule\noalign{}
\begin{minipage}[b]{\linewidth}\raggedright
Variable
\end{minipage} & \begin{minipage}[b]{\linewidth}\raggedright
Origen (NHANES)
\end{minipage} & \begin{minipage}[b]{\linewidth}\raggedright
Descripción
\end{minipage} & \begin{minipage}[b]{\linewidth}\raggedright
Codificación
\end{minipage} \\
\midrule\noalign{}
\endhead
\bottomrule\noalign{}
\endlastfoot
\texttt{edad} & RIDAGEYR & Edad al momento del examen & Años (80 = 80 o
más) \\
\texttt{sexo} & RIAGENDR & Sexo & Hombre / Mujer \\
\texttt{raza} & RIDRETH3 & Raza y origen hispano & 6 grupos \\
\texttt{educacion} & DMDEDUC2 & Nivel educativo & 5 niveles \\
\texttt{pir} & INDFMPIR & Razón ingreso familiar / línea de pobreza &
0-5 \\
\texttt{cafeina\_mg} & DR1TCAFF & Cafeína total ingerida en el día
(recordatorio 24h) & mg/día \\
\texttt{kcal} & DR1TKCAL & Energía total ingerida en el día &
kcal/día \\
\texttt{peso\_dieta} & WTDRD1 & Peso muestral del recordatorio dietario
& Numérico \\
\texttt{coronaria} & MCQ160C & Diagnóstico autorreportado de enfermedad
coronaria & 1 = sí, 0 = no \\
\texttt{infarto} & MCQ160E & Diagnóstico autorreportado de infarto & 1 =
sí, 0 = no \\
\texttt{acv} & MCQ160F & Diagnóstico autorreportado de ACV & 1 = sí, 0 =
no \\
\texttt{ecv} & derivada & Al menos uno de los tres diagnósticos
anteriores & 1 = sí, 0 = no \\
\texttt{cat\_cafeina} & derivada & Categoría de consumo de cafeína & 5
niveles \\
\texttt{grupo\_edad} & derivada & Grupo etario & 20-39, 40-59, 60+ \\
\end{longtable}
}

\section{Referencias}\label{referencias}

\protect\phantomsection\label{refs}
\begin{CSLReferences}{1}{1}
\bibitem[\citeproctext]{ref-cdc2018nhanes}
Centers for Disease Control and Prevention (CDC), National Center for
Health Statistics. 2018. \emph{National Health and Nutrition Examination
Survey Data 2017--2018}. U.S. Department of Health; Human Services.

\bibitem[\citeproctext]{ref-ding2014coffee}
Ding, Ming, Shilpa N. Bhupathiraju, Ambika Satija, Rob M. van Dam, y
Frank B. Hu. 2014. {«Long-term coffee consumption and risk of
cardiovascular disease: a systematic review and a dose-response
meta-analysis of prospective cohort studies»}. \emph{Circulation} 129
(6): 643-59.

\bibitem[\citeproctext]{ref-little2020statistical}
Little, Roderick J. A., y Donald B. Rubin. 2020. \emph{Statistical
Analysis with Missing Data}. 3rd ed. Wiley Series en Probability y
Statistics. John Wiley \& Sons.
\url{https://doi.org/10.1002/9781119482260}.

\bibitem[\citeproctext]{ref-poole2017coffee}
Poole, Robin, Oliver J. Kennedy, Paul Roderick, Jonathan A. Fallowfield,
Peter C. Hayes, y Julie Parkes. 2017. {«Coffee consumption and health:
umbrella review of meta-analyses of multiple health outcomes»}.
\emph{BMJ} 359: j5024.

\bibitem[\citeproctext]{ref-stevens2021coffee}
Stevens, Laura M., Erin Linstead, Jennifer L. Hall, y David P. Kao.
2021. {«Association between coffee intake and incident heart failure
risk: a machine learning analysis of the FHS, the ARIC Study, and the
CHS»}. \emph{Circulation: Heart Failure} 14 (2): e006799.

\end{CSLReferences}

\end{document}
